% ------------------------------------------------------------------------
% -*-TeX-*- -*-Hard-*- Smart Wrapping
% ------------------------------------------------------------------------


\chapter{結論}
\label{cha:7}

\section{本研究のまとめ}
\label{sec:7.1}

本研究は，性格の異なるペルソナを持つLoRAエージェントの社会が，議論を通じて問題を解き，議論から学び，チームへの貢献で次世代の代表を選ぶ進化的社会学習を提案し，事前に固定した計画のもとで検証した．予備研究の再検証で見つかった採点器の欠陥と，自分の回答を提示しない議論の実装を修正し，選抜には予備研究より大きな問題集合と全連合の実測を用いて推定誤差を減らす設計を採った．ペルソナは理論から網羅的に作った候補の中からデータと事前の規則で選んだ．

研究課題への答えは次のとおりである．

\begin{description}
\item[RQ1（進化の効果）] 系統Sの最終世代の社会の精度は，世代0の社会を上回らなかった（$-0.61$pt，$p = 0.15$）．差の90\%信頼区間は事前に定めたマージン$\pm 2$ptの内側にあり，両者は同等だった．3世代の進化的社会学習では，未見問題での精度は上がらなかった．% TODO(最終): seed 1〜4 の値に更新する．
\item[RQ2（選抜の効果）] チームへの貢献による選抜（系統S）と選抜なしの社会学習（系統N）の差は検出されず（$-0.04$pt，$p = 0.95$），両者は$\pm 2$ptの範囲で同等だった．
\item[RQ3（適応度の水準）] % TODO(最終): 系統A1の結果を記入する．
\item[RQ4（位置づけ）] 系統Sの最終世代の社会は，ベースモデル単体を$+3.13$pt，SC@3を$+1.66$pt有意に上回った．生成回数と出力トークン数を揃えたSC@4との差（$+1.10$pt，$p = 0.087$）は検出されず，SC@9およびRFTのSC@9とは$\pm 2$ptの範囲で同等だった．探索的には，世代0の社会がSC@4を$+1.71$pt上回り，SC@9の45\%の出力トークンでSC@9と同等の精度に達した．
\end{description}

以上から，本研究の主な知見は次の三つにまとめられる．第一に，4B級のモデルの3体の議論の社会は，ベースモデル単体とSC@3を有意に上回り，1問あたり約4回の生成でSC@9と同等の精度に達した．議論は議論なしの多数決を約1pt上回ったが，同程度の上積みは自分の回答を提示しない予備研究の手続きでも見られ，本研究のプロトコルに固有の効果とはいえない．計算量を揃えたSC@4を上回るかどうかは，事前登録した比較では確かめられなかった．第二に，同じベースモデルを共有するペルソナの社会では，議論の正解からの社会学習とチームへの貢献による選抜を重ねても，個体の精度，回答の多様性，社会の精度はいずれも上がらなかった．学習の信号が小さく，候補間の差が選抜の推定誤差を下回り，さらに親の値が上振れしたまま比べられたことが理由と考えられる．第三に，ペルソナの選び方は，理論と事前の規則とデータによって説明できるが，推論の手続きとしてのペルソナの違いが議論のチームの精度を左右するという証拠は得られなかった．

\section{今後の課題}
\label{sec:7.2}

本研究の結果は，進化的社会学習が機能するための条件を示唆している．今後の課題として次のものがある．

\begin{itemize}
\item \textbf{学習の信号を強めること}．自分がすでに解ける問題ではなく，チームとして初めて解けた問題や，議論で正解に転じた過程を重点的に学ぶ．世代あたりの学習の量を増やし，個体の精度が実際に上がる条件を先に確かめる．
\item \textbf{構成員の違いを作ること}．異なるベースモデルや異なる情報源を持つエージェントで社会を作り，他者から学ぶ価値のある違いを持たせる．人間の社会の文化的進化に近い状況である．
\item \textbf{選抜の推定誤差を下げること}．devを大きくし，逐次半減法のように有望な候補に評価を集中させて，候補間の小さな差を選べるようにする．親も各世代で新しい生成seedで測り直し，選抜の上振れを持ち越さないようにする．
\item \textbf{ボールドウィン型の対照}．学習の成果を子に受け渡さず，選抜の結果だけを受け継ぐ方式と比べ，ラマルク的な継承の効果を切り分ける．
\item \textbf{系統の反復}．同じ手続きを乱数を変えて複数回実行し，系統ごとのばらつきを含めて結論を確かめる．
\end{itemize}
